\subsection{阿佩尔多项式的生成级数}

设 E 为一个未定元 S 的形式级数环，其系数取在多项式环 K{[}X{]} 中（Alg., IV. 41）；换句话说，即形式级数 \(g(\mathrm{X}, \mathrm{S}) = \sum_{n=0}^{\infty} \alpha_n(\mathrm{X}) \mathrm{S}^n\) 组成的环，其中诸 \(\alpha_n\) 属于 K{[}X{]}。对每个算子 \(U\) （K{[}X{]} 上的），定义 E 到其自身的一个映射 \(U_{\mathrm{X}}\) 如下：置 \(U_{\mathrm{X}}(g(\mathrm{X}, \mathrm{S})) = \sum_{n=0}^{\infty} U(\alpha_n)\mathrm{S}^n\) 。显然，E 是以 K 为系数的未定元 S 的形式级数环 K{[}{[}S{]}{]} 上的一个模；由 \(U\) 在 K{[}X{]} 上的线性性，立即可以验证，对每个元素 \(\theta \in \mathrm{K}[[\mathrm{S}]]\) 和每个 \(g \in \mathrm{E}\) ，有 \(U_{\mathrm{X}}(\theta g) = \theta U_{\mathrm{X}}(g)\) ；换句话说，\(U_{\mathrm{X}}\) 是模 E 到其自身中的一个线性映射。

命题 6. 设 \(U = D^{p}V = u(D)\) 是 K{[}X{]} 上一个 p 阶复合算子，其中 \(u(S)\) 是 K{[}{[}S{]}{]} 中一个 p 阶形式幂级数。则

\[
U _ {\mathrm{X}} \big (\exp (\mathrm{XS}) \big) = u (\mathrm{S}) \exp (\mathrm{XS})\tag{15}
\]

\[
\frac {\mathrm{S} ^ {p}}{u (\mathrm{S})} \exp (\mathrm{XS}) = \sum_ {n = 0} ^ {\infty} \frac {1}{n !} u _ {n} (\mathrm{X}) \mathrm{S} ^ {n}\tag{16}
\]

其中 \(u_{n}\) 是与 U 相联系的指标为 n 的阿佩尔多项式。

由定理 1 的按语（VI, p.~271），为建立 (15)，只需证明对每个多项式 \(f(\mathrm{Y}) \in \mathrm{K}[\mathrm{Y}]\) 有

\[
U _ {\mathrm{X}} \left(\exp (\mathrm{XD} _ {\mathrm{Y}}) (f (\mathrm{Y}))\right) = u (\mathrm{D} _ {\mathrm{Y}}) \left(\exp (\mathrm{XD} _ {\mathrm{Y}}) (f (\mathrm{Y}))\right).\tag{17}
\]

而 (17) 中的第一项是 \(U_{\mathrm{X}}(f(\mathrm{X} + \mathrm{Y}))\) ，并且由于 \(U = u(\mathrm{D})\) ，(17) 中的第二项是 \(U_{\mathrm{Y}}(f(\mathrm{X} + \mathrm{Y}))\) ，于是恒等式 (17) 就归结为 (1) （VI, p.~270）。

现在只需把 (15) 应用于复合算子 \(V^{-1} = D^{p}/u(D)\) 即得 (16)，因为按定义有

\[
V ^ {- 1} \big (\exp (\mathrm{XS}) \big) = \sum_ {n = 0} ^ {\infty} \frac {1}{n !} u _ {n} (\mathrm{X}) \mathrm{S} ^ {n}.
\]

注意，把形式级数 \(S^{p}/u(S)\) 与 \(\exp(XS)\) 相乘并利用 (6)，也可以得到 (16)。

我们称形式级数 (16) 为与 U 相联系的阿佩尔多项式的生成级数。

